\subsection{Absolutely Semisimple Modules}

DEFINITION 1. --- Let K be a commutative field and A a K-algebra. We call an A-module M absolutely semisimple if the \(\mathrm{A}_{(\mathrm{L})}\) -module \(\mathrm{M}_{(\mathrm{L})}\) is semisimple for every extension L of K.

Every absolutely semisimple module is semisimple. Conversely, if the field K is perfect, then every semisimple A-module that is finite-dimensional over K is absolutely semisimple (VIII, p.~222, Corollary 1, a)) because every extension of a perfect field is separable (V, §7, No.~1, p.~36, Proposition 2).

PROPOSITION 1. --- Let K be a commutative field and A a K-algebra.

\begin{enumerate}
\def\labelenumi{\alph{enumi})}
\item
  Every direct sum of absolutely semisimple A-modules is an absolutely semisimple A-module. Every submodule and quotient module of an absolutely semisimple module is absolutely semisimple.
\item
  Let M be an A-module and let L be an extension of the field K. The A-module M is absolutely semisimple if and only if the \(\mathrm{A}_{(\mathrm{L})}\) -module \(\mathrm{M}_{(\mathrm{L})}\) is absolutely semisimple.
\end{enumerate}

Assertion a) follows from the analogous assertion for semisimple modules (VIII, p.~56, Corollaries 1 and 3).

Suppose that the A-module M is absolutely semisimple, and let \(L'\) be an extension of L. As an \(\mathrm{A}_{(\mathrm{L}')}\) -module, \(L' \otimes_{L} M_{(L)}\) is isomorphic to \(\mathrm{M}_{(\mathrm{L}')}\) (II, §5, No.~1, p.~273, Proposition 2); it is therefore a semisimple module. This proves that \(\mathrm{M}_{(\mathrm{L})}\) is absolutely semisimple.

Conversely, suppose that \(M_{(L)}\) is absolutely semisimple. Let \(L'\) be an extension of K. There exists a composite extension \((\Omega, u, v)\) of L and \(L'\) (V, §2, No.~4, p.~13, Corollary); we identify L and \(L'\) with subextensions of \(\Omega\) .

The \(A_{(\Omega)}\) -module \(M_{(\Omega)}\) is isomorphic to \((\mathrm{M}_{(\mathrm{L})})_{(\Omega)}\) ; it is therefore semisimple. But \(M_{(\Omega)}\) is also isomorphic to \((\mathrm{M}_{(\mathrm{L}^{\prime})})_{(\Omega)}\) , and Proposition 8, a) of VIII, p.~222 implies that \(M_{(L^{\prime})}\) is semisimple. So M is absolutely semisimple.

PROPOSITION 2. --- Let K be a commutative field, A a K-algebra, and M an A-module that is finite-dimensional over K. The following properties are equivalent:

\begin{enumerate}
\def\labelenumi{(\roman{enumi})}
\item
  The A-module M is absolutely semisimple.
\item
  There exists an extension \(\mathrm{P}\) of \(\mathrm{K}\) that is a perfect field such that the \(\mathrm{A}_{(\mathrm{P})}\) -module \(\mathrm{M}_{(\mathrm{P})}\) is semisimple.
\item
  The A-module M is semisimple, and the center of its commutant is an étale algebra over the field K.
\end{enumerate}

It is clear that (i) implies (ii). Assume that (ii) holds; then the \(A_{(P)}\) -module \(M_{(P)}\) is absolutely semisimple by Corollary 1, a) of VIII, p.~222. By Proposition 1, b), M is then absolutely semisimple. Therefore, (ii) implies (i).

Suppose that M is semisimple. Let Z be the center of the commutant of M; it is a commutative algebra of finite degree over the field K. If L is an extension of K, then it follows from Proposition 8, b) of VIII, p.~222 that the \(A_{(L)}\) -module \(M_{(L)}\) is semisimple if and only if the ring \(Z_{(L)}\) is reduced. The equivalence of (i) and (iii) therefore follows from Theorem 4 of V, §6, No.~7, p.~34.

